%%%%%%%%%%%%%%%%%%%%%%%%%%%%%%%%%%%%%%%%%
% Beamer Presentation
% LaTeX Template
% Version 1.0 (10/11/12)
%
% This template has been downloaded from:
% http://www.LaTeXTemplates.com
%
% License:
% CC BY-NC-SA 3.0 (http://creativecommons.org/licenses/by-nc-sa/3.0/)
%
%%%%%%%%%%%%%%%%%%%%%%%%%%%%%%%%%%%%%%%%%

%-------------------------------------------------------------------------------
%	PACKAGES AND THEMES
%-------------------------------------------------------------------------------

\documentclass{beamer}
\usepackage{xcolor}
\usepackage{graphicx}
\usepackage{tikz}
\usepackage{listings}
\usepackage{multicol}

\DeclareMathOperator{\diag}{diag}

\definecolor{applegreen}{rgb}{0.55, 0.71, 0.0}
\definecolor{blue(ncs)}{rgb}{0.0, 0.45, 0.60}
\definecolor{burgundy}{rgb}{0.5, 0.0, 0.13}

\definecolor{cadet}{rgb}{0.33, 0.41, 0.47}
\definecolor{airforceblue}{rgb}{0.36, 0.54, 0.66}

\mode<presentation> {

\usetheme{CambridgeUS}

\usecolortheme{wolverine}

\definecolor{gold}{HTML}{D4A017}
\definecolor{darkgold}{HTML}{B7950B}

\setbeamercolor{palette primary}{bg=cadet,fg=white}
\setbeamercolor{palette secondary}{bg=airforceblue,fg=white}
\setbeamercolor{palette tertiary}{bg=black,fg=white}
\setbeamercolor{palette quaternary}{bg=cadet,fg=white}

\setbeamercolor{frametitle}{bg=airforceblue,fg=white}

\setbeamercolor{section number projected}{bg=black,fg=cadet}
\setbeamercolor{item}{fg=black,bg=cadet}

\setbeamertemplate{page number in head/foot}[framenumber]

\lstset{basicstyle=\ttfamily\footnotesize,breaklines=true}
}

\lstdefinestyle{colouredC}{ 
  commentstyle=\color[gray]{0.4},
  keywordstyle=\bfseries,
  keywordstyle=[2]\color[rgb]{0.75, 0, 0},
  keywordstyle=[3]\color[rgb]{0, 0.5, 0},
  keywordstyle=[4]\color[rgb]{0, 0.5, 0},
  keywordstyle=[5]\color[rgb]{0, 0, 0.75},
  numberstyle=\tiny\color{mGray},
  stringstyle=\color[rgb]{0, 0, 0.5},
  basicstyle=\ttfamily,
  language=C,
  morekeywords=[3]{CeedInt, CeedScalar},
}

\lstdefinestyle{boxedC}{
  style=colouredC,
  numbers=left,
  firstnumber=auto,
  numberblanklines=true,
  frame=trbL,
  numberstyle=\tiny,
  frame=leftline,
  numbersep=7pt,
  framesep=5pt,
  framerule=10pt,
  xleftmargin=15pt,
  backgroundcolor=\color[gray]{0.97},
  rulecolor=\color[gray]{0.90}%
}

\usepackage{graphicx} % Allows including images
\usepackage{booktabs} % Allows the use of \toprule, \midrule and \bottomrule in tables

%----------------------------------------------------------------------------------------
%	TITLE PAGE
%----------------------------------------------------------------------------------------

\title[Ratel MI300A APUs]{Ratel - Using MI300A APUs with libCEED and PETSc} % The short title appears at the bottom of every slide, the full title is only on the title page

\author{Jeremy L Thompson} % Your name
\institute[CU Boulder] % Your institution as it will appear on the bottom of every slide, may be shorthand to save space
{University of Colorado Boulder \\ % Your institution for the title page
\medskip
\textit{jeremy@jeremylt.org} % Your email address
}
\date{11 Feb 2026} % Date, can be changed to a custom date

\begin{document}

\begin{frame}
\titlepage % Print the title page as the first slide
\end{frame}
 
%------------------------------------------------

\begin{frame}
\frametitle{Ratel Team}

\begin{center}
\includegraphics[height=2.75cm]{Ratellogo.png}
\end{center}

{\flushleft

Repository: https://gitlab.com/micromorph/ratel\\

~\\

Developers: Zach R. Atkins, Jed Brown, Fabio Di Gioacchino,\\
\hspace{19mm} Leila Ghaffari, Zach Irwin, Rezgar Shakeri,\\
\hspace{19mm} Ren Stengel, Jeremy L Thompson\\

The authors acknowledge support by the Department of Energy, National Nuclear Security Administration, Predictive Science Academic Alliance Program (PSAAP) under Award Number DE-NA0003962.

}

\begin{center}

\includegraphics[height=0.7cm]{psaap-center-logos}
\end{center}

\end{frame}

%------------------------------------------------

\begin{frame}
\begin{center}
\frametitle{Matrix-Free Operators from libCEED}

\includegraphics[height=5.0cm]{libCEEDAPI.png}

~\\

libCEED provides arbitrary order matrix-free FEM-like operators\\

\end{center}
\end{frame}

%------------------------------------------------

\begin{frame}[fragile]
\begin{center}
\frametitle{MatCEED}

\lstinline{MatCEED} wraps a linear \lstinline{CeedOperator}

{\tiny
\begin{lstlisting}[style=boxedC]
MatCeedCreate(DM dm_x, DM dm_y, CeedOperator op_mult, CeedOperator op_mult_transpose, Mat *mat);
MatCeedSetTime(Mat mat, PetscReal time);

// COO assembly support
MatCeedCreateMatCOO(Mat mat_ceed, Mat *mat_coo);
MatCeedSetPreallocationCOO(Mat mat_ceed, Mat mat_coo);
MatCeedAssembleCOO(Mat mat_ceed, Mat mat_coo);

// Private - Core functionality
MatMult_Ceed(Mat A, Vec X, Vec Y)
MatMultTranspose_Ceed(Mat A, Vec Y, Vec X)

// Private - PC support
MatGetDiagonal_Ceed(Mat A, Vec D);
MatGetDiagonalBlock_Ceed(Mat mat_ceed, Mat *mat_block);
MatInvertBlockDiagonal_Ceed(Mat mat_ceed, const PetscScalar **values);
MatInvertVariableBlockDiagonal_Ceed(Mat mat_ceed, PetscInt num_blocks, const PetscInt *block_sizes, PetscScalar *values);

\end{lstlisting}
}

~\\

\lstinline{MatCEED} can be used as a standard PETSc \lstinline{Mat}

\end{center}
\end{frame}

%------------------------------------------------

\begin{frame}
\begin{center}
\frametitle{MI300As}

\includegraphics[height=4.0cm]{UnifiedMemory.png}

~\\

Unified memory on MI300A removes CPU-GPU communication\\
(See also AI MAX+395)\\

~\\

Mildly slower memory but generally worth the tradeoff\\

\end{center}
\end{frame}

%------------------------------------------------

\begin{frame}
\begin{center}
\frametitle{Porting Outcomes}

\begin{itemize}

\item PETSc and libCEED changes to support APUs\\

~\\

\begin{itemize}

\item Allocate all memory as device memory\\

~\\

\item Treat same pointer as host memory\\

~\\

\item Must forbid mismatched memory operations (\lstinline{memset})\\

\end{itemize}

~\\

\item Unified memory improves performance on APU hardware\\

~\\

\begin{itemize}

\item At least 5\% speedup with unified memory\\

\end{itemize}

\end{itemize}

\end{center}
\end{frame}

%------------------------------------------------

\begin{frame}
\begin{center}
\frametitle{Ogden}

\includegraphics[height=6.5cm]{MPM_Ogden.png}\\

Can run MPM models efficiently\\
Reduced timestep setup costs (point migration, not shown)

\end{center}
\end{frame}

%------------------------------------------------

\begin{frame}
\begin{center}
\frametitle{Example - Press Simulation}

\includegraphics[height=5cm]{iMPMCloseup.png}
\hspace*{1cm}
\includegraphics[height=5cm]{iMPMCompress.png}

~\\

Compression of mock HE grains (gold) and binder (pink) mixture\\

(Reset background mesh to computational region on each timestep)\\

Latest runs use $\sim$ \hspace{-6pt} 10B material points

\end{center}
\end{frame}
 
%------------------------------------------------

\begin{frame}
\frametitle{Questions?}

\begin{center}
\includegraphics[height=2.75cm]{Ratellogo.png}
\end{center}

{\flushleft

Repository: https://gitlab.com/micromorph/ratel\\

~\\

Ratel Team: Zach R. Atkins, Jed Brown, Fabio Di Gioacchino,\\
\hspace{20mm} Leila Ghaffari, Zach Irwin, Rezgar Shakeri,\\
\hspace{20mm} Ren Stengel, Jeremy L Thompson\\

~\\

Grant: Predictive Science Academic Alliance Program (DE-NA0003962)\\

}

\begin{center}

\includegraphics[height=0.8cm]{psaap-center-logos}
\end{center}

\end{frame}

%-------------------------------------------------------------------------------

\end{document}

%-------------------------------------------------------------------------------
